\emph{Översikt}
Åldersprogrammet i föreläsningen \emph{Variabler och utskrifter} frågade
den som körde det efter namn och födelseår, men vi sa bara ett par
meningar om hur det gick till.  Här tar vi upp de raderna igen.  Vi
använder \mintinline{python}{input} för att läsa in text, och upptäcker
att det som kommer tillbaka alltid är en \emph{sträng} --- även när
användaren skriver siffror.  Därför behöver vi \emph{typkonvertering},
och därför behöver vi veta vad datatypen gör med en operation:
\mintinline{python}|"3" + "4"| och \mintinline{python}{3 + 4} ser lika
ut men ger olika svar.  Vi går igenom de vanligaste operationerna på tal
och strängar, och ett par strängmetoder som vi får användning för när vi
ska granska det användaren skrev, och vi utformar frågorna så att den
som kör programmet förstår vad som förväntas.  Sedan svarar användaren
något vi inte kan räkna med, och programmet går sönder: körningen
\emph{avbryts} och Python skriver en \emph{spårutskrift}
(\foreignlanguage{english}{traceback}).  Vi lär oss läsa den, och sedan
ta hand om felet själva med \mintinline{python}{try} och
\mintinline{python}{except}: en gren per sorts fel, de specifika först
och det allmänna sist eller inte alls, och med ett meddelande skrivet
för användaren i stället för för programmeraren.  Och till sist ser vi
hur våra egna funktioner kan \emph{kasta} särfall vidare till den som
anropat dem.

\emph{Lärandemål}
Efter denna modul ska studenten kunna:
\begin{restatable}{lo}{InputLORead}\label{InputLORead}%
  Läsa in text från användaren med \mintinline{python}{input} och
  förklara att resultatet alltid är en sträng, oavsett vad användaren
  skrev.
\end{restatable}% Lo08
\begin{restatable}{lo}{InputLOConvert}\label{InputLOConvert}%
  Omvandla en inläst sträng till \mintinline{python}{int} eller
  \mintinline{python}{float} med typkonvertering, och avgöra var i
  programmet omvandlingen hör hemma.
\end{restatable}% Lo08, Lo03
\begin{restatable}{lo}{InputLOTypes}\label{InputLOTypes}%
  Använda de vanligaste operationerna och metoderna på tal och strängar,
  och förutsäga vilken typ och vilket värde resultatet får.
\end{restatable}% Lo03
\begin{restatable}{lo}{InputLOPrompt}\label{InputLOPrompt}%
  Utforma frågor, utskrifter och felmeddelanden så att den som kör
  programmet förstår vad som förväntas, vad svaret betydde och vad hon
  kan göra när något gick fel --- i stället för att låta programmet
  krascha eller vidarebefordra Pythons egen text.
\end{restatable}% Lo05 (användbara utskrifter), Lo08 (interaktiva program)
\begin{restatable}{lo}{ExcLOException}\label{ExcLOException}%
  Förklara vad ett särfall (\foreignlanguage{english}{exception}) är,
  och redogöra för vad som händer när ett särfall inträffar.
\end{restatable}% Lo04 (styrstrukturer), Lo08 (interaktiva program)
\begin{restatable}{lo}{ExcLOTraceback}\label{ExcLOTraceback}%
  Läsa en spårutskrift (\foreignlanguage{english}{traceback}) och ur
  den utläsa var felet uppstod, vilken sorts fel det var och vad som var
  fel.
\end{restatable}% Lo08 (interaktiva program), Lo11 (granska program)
\begin{restatable}{lo}{ExcLOCatch}\label{ExcLOCatch}%
  Fånga särfall med \mintinline{python}{try} och
  \mintinline{python}{except} --- de specifika sorterna före den
  allmänna --- och redogöra för vad som händer i programmet.
\end{restatable}% Lo04 (styrstrukturer), Lo08 (interaktiva program)
\begin{restatable}{lo}{ExcLORaise}\label{ExcLORaise}%
  Läsa och skriva \mintinline{python}{raise} i en egen funktion, och
  förklara hur ett särfall sprider sig uppåt till den som anropat
  funktionen.
\end{restatable}% Lo03 (dela upp program), Lo04 (styrstrukturer)

\emph{Förkunskaper}
Studenten bör ha tagit del av föreläsningen \emph{Funktioner}:
variabler, datatyperna \mintinline{python}{int}, \mintinline{python}{float}
och \mintinline{python}{str}, utskrifter med \mintinline{python}{print}
och egna funktioner med parametrar och returvärden.

\emph{Läsning}
Kompletterande material finns i kursens videogenomgång \emph{Inmatning
och felhantering}.  Pythons egen dokumentation beskriver särfallen och
hur man fångar dem i kapitlet
\emph{Errors and Exceptions}.\footnote{\label{fn:pydoc}%
  \url{https://docs.python.org/3/tutorial/errors.html} (hämtad
  2026-09-03).}
